\section{Architecture Design and Training Stability}
\label{sec:architecture_stability}

\begin{figure}[htbp]
    \centering
    \begin{subfigure}[t]{0.48\linewidth}
        \centering
        \includegraphics[width=\linewidth]{architecture/2B_v3.pdf}
        \caption{Activation statistics of the baseline architecture}
        \label{fig:2B_v3}
    \end{subfigure}
    \hfill
    \begin{subfigure}[t]{0.48\linewidth}
        \centering
        \includegraphics[width=\linewidth]{architecture/2B_v4.pdf}
        \caption{Activation statistics after applying Sandwich Normalization and Logits Soft-capping}
        \label{fig:2B_v4}
    \end{subfigure}
    
    \caption{Comparison of internal activation magnitudes before and after architectural optimization. The experiment is conducted with a 3B model.}
    \label{fig:main_figure_label}
\end{figure}

\sys is trained on Huawei Ascend 910A accelerators, which are similar to NVIDIA V100s in supporting only FP16 precision. However, FP16 has a limited dynamic numerical range, which introduces overflow risks when model parameters or activations grow too large. To keep training stable, we first identify the activations that are most likely to overflow and then introduce structural changes that keep their values within safe bounds.

Following the standard Llama architecture, the model uses SwiGLU~\cite{dauphin2017swiglu}, RMSNorm~\cite{zhang2019rmsnorm}, and RoPE~\cite{su2024rope}. We adopt mixed precision training, where operators that need higher precision, such as Softmax and RMSNorm, run in FP32, and the remaining computations run in FP16. Despite this setup, training on large and diverse datasets, including code and mathematics, still leads to strong numerical instability. As shown in \Cref{fig:2B_v3}, most instability comes from two places: the attention logits and the activations after the SwiGLU function in the MLP layers. In practice, the maximum activation values grow without control. They exceed $10{,}000$ after processing one trillion tokens, which is close to the FP16 upper limit. As a result, the dynamic loss scaler decreases its scaling factor to avoid overflow. This drop pushes many gradients below the FP16 minimum representable value, which causes underflow. The gradients then become inaccurate, harming convergence and sometimes causing training to fail.

To solve these issues, we use Logits Soft-Capping~\cite{bello2017softcapping} and Sandwich Normalization~\cite{ding2021cogview}. This follows the design choices of Gemma 2~\cite{morgane2024gemma2}. These techniques place strict bounds on activation values. As shown in \Cref{fig:2B_v4}, soft-capping reduces the L1 norm of attention logits by about an order of magnitude. At the same time, sandwich normalization reduces the accumulation of large values in residual connections and keeps the L1 norm of MLP activations within a safe range. To further improve stability, we set the weight decay to $0.1$, apply soft-capping to the final output logits, and replace the soft-capping inside each attention layer with QK-Norm~\cite{henry2020qknorm}. The full configuration of \sys is listed in \Cref{tab:training_config} and the implementation details are discussed in \Cref{sec:implement_architecture}.







